% Capitolul 2: stadiul actual (TH-2).
% Fiecare citare trimite la o intrare din referinte.bib verificată la sursă;
% afirmațiile despre proiectul acesta trimit la capitolele lui, nu la citări.

\chapter{Stadiul actual}
\label{cap:stadiu}

Un cititor de cărți electronice răspunde la trei întrebări: pe ce se afișează
textul, ce procesor îl pregătește și în ce formă ajunge cartea la acel
procesor. Capitolul trece prin răspunsurile date până acum, de la tehnologia
ecranului la cititoarele comerciale și la proiectele deschise construite pe
microcontrolere, apoi la formatele de fișier și la algoritmii de rupere a
textului în rânduri. La final arată locul acestei lucrări printre ele.

% ---------------------------------------------------------------------------
\section{Hârtia electronică}
\label{sec:stadiu-epaper}

Ecranele cititoarelor de azi descind din cerneala electroforetică descrisă de
Comiskey și colaboratorii săi în 1998 \citep{comiskey1998}: particule încărcate
electric, închise în microcapsule, care se deplasează într-un câmp continuu
spre fața sau spre spatele celulei. Lucrarea numește proprietățile care fac
tehnologia potrivită pentru citit. Imaginea e bistabilă: rămâne și după ce
tensiunea dispare. Adresarea cere puțină energie. Ecranul e reflectiv, cu un
contrast și un aspect apropiate de cerneala pe hârtie. Microîncapsularea
rezolvă problemele de durată de viață ale variantelor anterioare, iar ecranul
se poate fabrica în întregime prin imprimare.

Bistabilitatea are o consecință directă pentru un cititor: energia se consumă
doar când se schimbă pagina, nu cât timp pagina stă pe ecran. Prețul e viteza.
Particulele au nevoie de sute de milisecunde ca să ajungă în poziție, iar
tensiunea nu se aplică într-un singur impuls, ci ca o formă de undă cu mai
multe faze, care întâi activează particulele și abia apoi le duce în starea
dorită \citep{he2020}. Formele de undă sunt proprii fiecărui panou, deci le
definește producătorul. O trecere în revistă din 2011 a
tehnologiilor de hârtie electronică \citep{heikenfeld2011} le compară pe toate
cele importante, arată limitele care vin din fizica dispozitivelor și
conchide că între ecranele reflective și tiparul pe hârtie rămâne o distanță
semnificativă.

Există două feluri de a comanda un astfel de panou. Panourile
cititoarelor comerciale au o interfață paralelă: un controler generează
pentru fiecare cadru tensiunile de rând și de coloană după forma de undă.
Procesoarele cititoarelor au controlerul integrat. Descrierea hardware a
SoC-ului i.MX~6SoloLite din nucleul Linux, de exemplu, conține un bloc EPDC
(\emph{electrophoretic display controller}) cu ceasurile lui
\citep{linux-imx6sl}. Proiectul deschis EPDiy face același lucru cu un ESP32:
o placă de comandă pentru panourile paralele recuperate din cititoare vechi,
cu 16 nuanțe de gri, plecând de la observația că metoda de comandă e aproape
aceeași pentru toate panourile matriceale e-ink \citep{epdiy}.

Panourile mici vin, în schimb, cu controlerul în modul. Panoul folosit aici,
E037A75 de 3,7~inci, are 416~×~240 de pixeli și un circuit UC8253 care
conține bufferele de rând și de coloană, temporizarea, convertorul DC-DC,
memoria cadrului și tabelele formelor de undă. Formele de undă sunt scrise de
producător în memoria OTP, iar comunicarea se face prin SPI cu trei sau patru
fire \citep{e037a75}. Microcontrolerul trimite deci doar imaginea paginii;
tensiunile și formele de undă nu sunt treaba lui.

% ---------------------------------------------------------------------------
\section{Cititoarele comerciale}
\label{sec:stadiu-comerciale}

Cititoarele Kindle, Kobo și PocketBook sunt calculatoare mici cu Linux. Kobo
publică sursele nucleului Linux și ale încărcătorului U-Boot folosite pe
cititoarele sale, împreună cu construcția bibliotecii Qt pentru dispozitive
încorporate pe care rulează interfața \citep{kobo-gpl}. Lanțurile de
compilare încrucișată întreținute de proiectul KOReader au câte o țintă ARM
Linux pentru fiecare familie: Kindle, Kobo, PocketBook, reMarkable și altele
\citep{koxtoolchain}. Dezasamblarea unui Kindle Voyage arată componentele unui
astfel de dispozitiv: un procesor Freescale i.MX~6SoloLite, memorie LPDDR2
externă, stocare eMMC și un circuit dedicat de gestionare a alimentării
\citep{ifixit-voyage}. Același i.MX~6SoloLite are, cum s-a văzut mai sus,
controlerul de hârtie electronică integrat.

Cu un sistem de operare complet și memorie externă, cartea se poate deschide
în formatul ei original, cum face KOReader pe aceleași dispozitive
(secțiunea~\ref{sec:stadiu-deschise}). Arhitectura e potrivită pentru un
produs comercial, dar e departe de ce se poate construi și înțelege în
întregime ca proiect de diplomă: procesorul, memoria externă și sistemul de
operare aduc complexitate, iar platformele sunt închise, argumentul de la care
pornește The Open Book.

% ---------------------------------------------------------------------------
\section{Proiecte deschise}
\label{sec:stadiu-deschise}

\paragraph{Programe pentru cititoarele existente.} KOReader e o aplicație de
citit cu licență AGPL-3.0 care rulează pe cititoarele Cervantes, Kindle, Kobo,
PocketBook și reMarkable, pe Android și pe calculatoare cu Linux, și deschide
PDF, DjVu, EPUB, FB2 și alte formate \citep{koreader}. E dovada că partea de
program a unui cititor poate fi deschisă. Rămâne însă legată de hardware-ul
comercial și de Linux.

\paragraph{Plăci cu ecran e-paper.} Inkplate e o familie de plăci cu ESP32,
Wi-Fi și Bluetooth, însoțite de o bibliotecă Arduino \citep{inkplate}. Ca și
EPDiy, oferă ecranul și comanda lui, nu un cititor: aranjarea textului rămâne
în seama aplicației.

\paragraph{Cititoare pe microcontrolere.} Proiectele care construiesc
cititorul întreg pe un microcontroler au o resursă rară: memoria. Un ESP32-S3
are 512~KB de SRAM internă \citep{esp32s3}. Trei proiecte arată trei răspunsuri
la această limită.

\emph{The Open Book} \citep{openbook} a pornit în 2019 cu ideea unui
dispozitiv de citit deschis, ca alternativă la platformele închise ale
cititoarelor populare. A fost refăcut în jurul plăcii Raspberry Pi Pico, cu
baterii AAA și cu fișierele de asamblare pentru JLCPCB publicate. Firmware-ul
lui, libros, nu citește EPUB: formatul lui canonic e textul simplu, cu titlul
pe primul rând și, opțional, un antet și coduri de control ASCII pentru
capitole și formatare. Cartea e simplificată înainte să ajungă pe dispozitiv.
Rândurile se rup lacom: se ia ultima ocazie de rupere care mai încape, fără
despărțire în silabe \citep{libros}.

\emph{diy-esp32-epub-reader} \citep{atomic14} face opusul: deschide EPUB direct
pe ESP32, cu miniz pentru arhivă și TinyXML2 pentru XHTML, și ignoră CSS. Cere
PSRAM, pentru că, după README-ul proiectului, parsarea fișierelor EPUB are
nevoie de multă memorie. Paragrafele se rup în rânduri prin programare
dinamică la nivel de cuvânt, minimizând suma pătratelor spațiului rămas la
capătul rândurilor, fără despărțire în silabe. Ultima modificare e din 2022.

\emph{CrossPoint} \citep{crosspoint} e cel mai nou și cel mai activ dintre ele. E un firmware
pentru cititorul Xteink X4, cu ESP32-C3, și pentru alte câteva dispozitive cu
ESP32-S3. Citește EPUB 2 și 3 pe dispozitiv, cu kerning și despărțire în
silabe. README-ul proiectului spune că ESP32-C3 are doar circa 380~KB de RAM
utilizabilă. Răspunsul e un cache agresiv pe card: fiecare carte deschisă
primește un dosar cu metadatele, regulile CSS și câte un fișier
cu așezarea fiecărui capitol. Codul de rupere în rânduri are două ramuri. Fără
despărțire în silabe, aplică aceeași programare dinamică cu pătratul spațiului
rămas. Cu despărțire, trece la un algoritm lacom care desparte doar cuvântul
care nu mai încape pe rând. CrossPoint mai citește un al doilea format, XTC,
care stochează fiecare pagină ca imagine gata randată de un convertor pe PC.
Pentru dispozitiv e cel mai ieftin format posibil, dar litera și așezarea se
fixează la conversie.

Tabelul~\ref{tab:stadiu-cititoare} rezumă cum tratează aceste proiecte
formatul și ruperea în rânduri, alături de proiectul de față.

\begin{table}[htbp]
  \centering
  \small
  \caption{Cititoare deschise pe microcontrolere: formatul citit pe dispozitiv
    și ruperea în rânduri.}
  \label{tab:stadiu-cititoare}
  \begin{tabularx}{\textwidth}{@{}>{\raggedright\arraybackslash}p{0.21\textwidth}
      >{\raggedright\arraybackslash}p{0.17\textwidth}
      >{\raggedright\arraybackslash}X >{\raggedright\arraybackslash}X@{}}
    \toprule
    Proiect & Procesor & Formatul pe dispozitiv & Ruperea în rânduri \\
    \midrule
    The Open Book (libros) & Raspberry Pi Pico & text simplu cu coduri de control
      & lacomă, fără despărțire \\
    diy-esp32-epub-reader & ESP32 cu PSRAM & EPUB & optimă, fără despărțire \\
    CrossPoint & ESP32-C3, ESP32-S3 & EPUB; XTC (pagini imagine)
      & optimă fără despărțire sau lacomă cu despărțire \\
    proiectul de față & ESP32-S3 & IBV (text structurat, cratime calculate)
      & Knuth–Plass, cu despărțire \\
    \bottomrule
  \end{tabularx}
\end{table}

% ---------------------------------------------------------------------------
\section{Formatele de fișier}
\label{sec:stadiu-formate}

EPUB, standardizat acum de W3C \citep{epub33}, e formatul obișnuit al cărților
electronice. O publicație EPUB e un set de documente XHTML, foi de stil CSS,
imagini și fonturi, descrise de un document de pachet și strânse într-o arhivă
ZIP (containerul OCF) cu extensia \texttt{.epub}. Intrările arhivei sunt fie
stocate necomprimat, fie comprimate cu Deflate \citep{epub33}, algoritmul
descris în RFC~1951 \citep{rfc1951}. Ca să afișeze o pagină dintr-un EPUB,
un dispozitiv trebuie deci să decomprime un document din arhivă, să parseze
XML, să aplice măcar o parte din CSS și abia apoi să așeze textul. Pe un
microcontroler asta înseamnă fie memorie externă (PSRAM, ca la
diy-esp32-epub-reader), fie un cache pe card refăcut la fiecare carte nouă
(ca la CrossPoint).

FictionBook~2 (FB2) e un format XML mai simplu: o carte e un singur fișier,
cu structura (secțiuni, titluri, paragrafe, citate, poezii cu strofe și
versuri) descrisă de o schemă XSD \citep{fb2}. Structura e mai apropiată de ce are
nevoie un cititor, dar fișierul rămâne XML de parsat pe dispozitiv.

La celălalt capăt sunt formatele pregătite pe PC: textul simplu al lui libros
\citep{openbook}, care pierde majoritatea structurii, și paginile imagine ale
formatului XTC \citep{crosspoint}, care păstrează aspectul dar nu mai permit
schimbarea literei. Între ele rămâne loc pentru un format intermediar: o carte
convertită pe PC, cu structura și despărțirea în silabe deja calculate, dar
cu așezarea în pagină lăsată dispozitivului. Acesta e formatul IBV, descris în
capitolul despre format.

% ---------------------------------------------------------------------------
\section{Algoritmii de rupere în rânduri}
\label{sec:stadiu-rupere}

Ruperea unui paragraf în rânduri are doi pași. Primul găsește locurile unde
textul se poate rupe: spații, cratime, alte semne. Anexa~14 a standardului
Unicode \citep{uax14} definește acest pas pentru toate scrierile și spune
explicit că alegerea rupturilor dintre aceste ocazii, ce înseamnă o alegere
optimă și alinierea rândurilor la ambele margini nu țin de ea, ci de
programele de nivel mai înalt care cunosc lățimea rândului. Al doilea pas
alege rupturile, iar aici algoritmii diferă.

Algoritmul lacom (\emph{first fit}) pune pe fiecare rând cât de mult text
încape și trece la următorul. E simplu și rapid, dar fiecare rând se decide
fără să se țină seama de cele care urmează. Un rând bine umplut poate lăsa
următorului un cuvânt lung care nu încape, iar rândul acela iese foarte rar,
cu spații mari între cuvinte.

Knuth și Plass \citep{knuth1981} au propus în 1981 ca paragraful să fie
tratat ca un întreg, astfel încât aspectul unui rând să poată fi influențat de
textul de pe rândurile următoare. Paragraful se descrie cu trei primitive:
cutii (literele), spații elastice (\emph{glue}), care se pot lăți sau îngusta,
și penalizări, care marchează locurile unde ruperea are un cost. Un singur
algoritm găsește rupturile optime, fără reveniri, prin programare dinamică.
E algoritmul din TeX. Hirschberg și Larmore \citep{hirschberg1987} au
formulat problema generală din spatele lui, subșirul de cost minim, au arătat
că un caz particular se rezolvă în $O(n \log n)$, iar pentru unele funcții de
cost în timp liniar, și au dat formarea optimă a paragrafului ca aplicație.

Minimizarea sumei pătratelor spațiului rămas, folosită de
diy-esp32-epub-reader și de CrossPoint, e un caz mai simplu din aceeași
familie: costul unui rând e dat doar de spațiul rămas gol la capătul lui, nu
de cât se lățesc spațiile la justificare. Modelul Knuth–Plass merge mai departe: costul unui
rând justificat depinde de cât se lățesc sau se îngustează spațiile, iar
despărțirea în silabe intră ca penalizare. Algoritmul hotărăște singur dacă o
cratimă merită pentru a evita un rând prea rar.

Pentru despărțirea în silabe, metoda folosită de TeX e cea a lui Liang
\citep{liang1983}, dezvoltată pentru TeX82. Problema e tratată ca una de
compresie a datelor: în loc de un dicționar cu toate cuvintele despărțite, se
păstrează un set mic de modele care permit sau interzic despărțirea între
anumite litere. Modelele pentru zeci de limbi, inclusiv română și rusă, sunt
întreținute în colecția hyph-utf8 \citep{texhyphen}.

Navigatoarele web au început să meargă în aceeași direcție. Ciorna CSS Text
Level~4 \citep{csstext4} definește \texttt{text-wrap-style: pretty}, cu care
navigatorul ar trebui să prefere calitatea vitezei și să ia în calcul mai
multe rânduri când decide unde rupe. Specificația lasă însă algoritmul exact
la latitudinea fiecărui navigator.

% ---------------------------------------------------------------------------
\section{Locul acestei lucrări}
\label{sec:stadiu-loc}

Proiectele studiate se împart după locul unde e făcută munca grea. Cititoarele
comerciale și KOReader au procesoare cu Linux și memorie externă, așa că
deschid EPUB direct. Cititoarele pe microcontrolere fie fac la fel cu memorie
externă sau cache pe card (diy-esp32-epub-reader, CrossPoint), fie mută
pregătirea cărții pe PC și pierd ceva la conversie: structura, în textul
simplu al lui libros, sau posibilitatea de a schimba litera, în paginile
imagine XTC.

Lucrarea de față ia o cale între ele. Un convertor pe PC citește EPUB, FB2 și
text simplu și produce un format propriu, IBV. Formatul păstrează structura
cărții (paragrafe, titluri, citate, versuri) și cratimele moi calculate cu
modelele lui Liang, iar textul e împărțit în blocuri comprimate cu Deflate,
fiecare de cel mult 32~KB. Formatul e gândit ca dispozitivul, un ESP32-S3,
să nu aibă nevoie de memorie externă: decomprimă câte un bloc și așază textul cu fontul și mărimea
alese de cititor. Paragrafele justificate se rup cu încadrarea totală
Knuth–Plass, cu parametrii din TeX. Ruperile sunt comparate, pe tot corpusul
de test, cu cele ale lui LuaTeX.

Dintre cititoarele deschise pe microcontrolere analizate aici, niciunul nu
combină ruperea optimă a paragrafului cu despărțirea în silabe: libros rupe
lacom, diy-esp32-epub-reader nu desparte în silabe, iar CrossPoint alege fie
optimizarea, fie despărțirea, nu amândouă. Panoul, cu controlerul și formele
de undă integrate, scutește microcontrolerul de comanda propriu-zisă a
ecranului. Rămân de rezolvat așezarea textului, formatul și consumul, care
sunt subiectul capitolelor următoare.
